\documentclass[11pt,a4paper]{article}

\usepackage[T1]{fontenc}
\usepackage[utf8]{inputenc}
\usepackage{lmodern}
\usepackage[margin=2.4cm]{geometry}
\usepackage{booktabs}
\usepackage{array}
\usepackage{enumitem}
\usepackage{xcolor}
\usepackage{siunitx}
\usepackage[colorlinks=true, linkcolor=blue!60!black, urlcolor=blue!60!black]{hyperref}
\usepackage{parskip}

\newcommand{\code}[1]{\texttt{#1}}
\newcolumntype{L}[1]{>{\raggedright\arraybackslash}p{#1}}

\title{Conference Roadmap --- P2 Basket Pad-Micromegas:\\
       from SPS Beam Test to Cosmic-Bench Characterization}
\author{Alexandra Kallitsopoulou}
\date{Draft roadmap, 11 July 2026 --- anchor the timeline to the actual conference date once fixed}

\begin{document}

\maketitle

%=====================================================================
\section{The story to tell}
%=====================================================================

One sentence: \emph{``We are developing pad-Micromegas detectors for high-rate
environments; we validated two geometries and optimized the VMM3a front-end at
the SPS, we are now characterizing the next detector generation on a cosmic
bench with full tracking, and we return to beam with VMM to close the loop.''}

The talk has a natural three-act structure.

\subsection*{Act 1 --- Beam test at SPS (Nov.\ 2025): geometry + electronics optimization}

\begin{itemize}[itemsep=2pt]
  \item Two detector geometries under test simultaneously with VMM3a readout:
    \begin{itemize}
      \item \textbf{P2 Large detector} --- fan-pad geometry, 256 pads (VMMs 12--15)
      \item \textbf{P2 Small detectors 1 \& 3} --- compact pad geometry (VMMs 8--11)
    \end{itemize}
  \item \textbf{SNR optimization}: full VMM3a shaping scan
    (gain $sg = 3.0/4.5/6.0$~mV/fC $\times$ peaking time $snt = 50/100/200$~ns),
    signal/noise run pairs via neighbour-triggering, MAD noise baseline + MPV
    signal $\rightarrow$ SNR per VMM and per channel.
    $\Rightarrow$ \emph{``We found the optimal operating point of the electronics.''}
  \item \textbf{Rate performance}: identical config points taken at
    ${\sim}5$~kHz and ${\sim}15$~kHz muon rates, plus \textbf{pion runs} with
    localized, much higher rate density.
    $\Rightarrow$ \emph{``We probed the detector where it will actually have to live.''}
  \item Trigger-based efficiency with accidental subtraction ($\Delta t$ peak
    fit + sideband) and per-pad efficiency maps on the real Gerber geometry.
\end{itemize}

\subsection*{Act 2 --- Cosmic bench (2026): new detector generation, precision characterization}

\begin{itemize}[itemsep=2pt]
  \item New P2 detectors (P2\_1 ``det1'', P2\_2 ``det2''), 256-pad fan
    geometry, read out with DREAM FEUs behind the \textbf{M3 beam telescope}
    $\rightarrow$ track-projected, \emph{geometry-resolved} measurements a beam
    trigger alone cannot give:
    \begin{itemize}
      \item pad-level and sliding-window \textbf{efficiency maps} on the true pad footprint
      \item \textbf{HV turn-on curve} (mesh 345--420~V, $\varepsilon$: $0.40 \rightarrow 0.66$ measured so far)
      \item \textbf{timing} from waveform TOA (${\sim}46$~ns vs.\ trigger, drift-geometry limited)
      \item \textbf{stability}: HV spark monitoring/veto, per-pad spark
        flagging, pedestal QA, long-run duty-cycle tracking (this is how the
        det1 intermittent-gain problem was \emph{found} --- that is a selling
        point: the bench methodology catches real detector pathologies)
    \end{itemize}
  \item $\Rightarrow$ \emph{``Before the next beam test, every detector is fully
    mapped, efficiency-calibrated and stability-vetted on the bench.''}
\end{itemize}

\subsection*{Act 3 --- Outlook: next beam test with VMM}

\begin{itemize}[itemsep=2pt]
  \item Bench-validated detectors + SPS-optimized VMM configuration
    $\rightarrow$ next campaign measures the thing that matters:
    \textbf{efficiency and stability vs.\ particle rate at the optimal working
    point}, with known geometry and known electronics response.
  \item The cross-electronics angle is a strength, not a complication:
    \textbf{DREAM (full waveforms, external tracking) for understanding, VMM3a
    (self-triggered, high rate) for operating.} Same detectors, two
    complementary readouts, two beam environments.
\end{itemize}

%=====================================================================
\section{What already exists (asset inventory)}
%=====================================================================

\begin{table}[h!]
\centering\small
\begin{tabular}{L{5.2cm} L{5.6cm} L{4.4cm}}
\toprule
\textbf{Analysis} & \textbf{Code} & \textbf{State} \\
\midrule
SNR config-scan pipeline (VMM + channel level) &
  \code{SNR\_Analysis/vmm\_config\_scan\_analysis.py} + modules &
  \textbf{Done} for 5 \& 15~kHz; report \code{vmm\_analysis\_report.html} \\
\addlinespace
Trigger rate / spill structure / occupancy &
  \code{vmm\_trigger\_analysis.py} + \code{trigger\_rate\_analysis\_report.md} &
  \textbf{Done} for muon datasets \\
\addlinespace
Trigger-coincidence efficiency + per-pad maps &
  \code{vmm\_detector\_efficiency.py} (FanPadDetector, validated REVERSE ordering) &
  Works on single runs; needs a config-scan / rate sweep \\
\addlinespace
Pad geometry models (Gerber-true) &
  \code{FanPadDetector.py}, \code{SquarePadDetector.py}, \code{Detector\_Mapping/} &
  \textbf{Done}, M3-cross-validated \\
\addlinespace
Cosmic-bench full QA pipeline (stages 02--13) &
  \code{cosmic\_bench\_analysis/} + \code{run\_p2\_pipeline.sh} &
  \textbf{Done} for det1 (4 long runs), det2 (1 run); PDF builders exist \\
\addlinespace
HV-scan efficiency &
  stage 11 + \code{build\_hv\_scan\_pdf.py} &
  Done for det1 345--420~V (365~V file to re-fetch) \\
\addlinespace
Timing (waveform TOA) &
  stage 13 &
  Working on det2, \textbf{not committed yet} \\
\addlinespace
Structural bias audit &
  stage 12 &
  Done (knob plateaus, estimator checks) --- backs up systematic error bars \\
\bottomrule
\end{tabular}
\end{table}

\textbf{Missing / to build:} pion-run analysis (nothing in the codebase
mentions pions), efficiency-vs-rate sweep across the muon datasets,
cross-geometry comparison plots, det2 HV scan, and the final cross-campaign
summary figures.

%=====================================================================
\section{Analysis roadmap (work packages, in priority order)}
%=====================================================================

\subsection*{WP1 --- SNR result finalization \emph{(mostly done --- polish only)}}

\begin{enumerate}[itemsep=1pt]
  \item Freeze the SNR summary: \textbf{SNR vs.\ $(sg, snt)$} matrix per
    detector, 5~kHz and 15~kHz side by side.
  \item One channel-level SNR map per geometry at the best config (shows
    uniformity, not just averages).
  \item Extract the headline number: \emph{best config $= (sg, snt)$ with
    SNR $= X$ on large, $Y$ on small.}
\end{enumerate}

\textbf{Hero plots:} SNR heat matrix ($sg \times snt$) per detector; example
signal-vs-noise ADC spectra at best/worst config (the ``why this matters''
plot); channel-level SNR pad map.

\subsection*{WP2 --- Rate performance, muons \emph{(1--2 weeks of work; reuses existing pieces)}}

\begin{enumerate}[itemsep=1pt]
  \item Run \code{vmm\_detector\_efficiency.py} over the matched config points
    of the 5~kHz \textbf{and} 15~kHz run tables $\rightarrow$ efficiency per
    detector per rate.
  \item Occupancy/hit-rate per pad on-spill (from trigger-analysis machinery)
    $\rightarrow$ rate-density maps.
  \item Check for rate-dependent effects: $\Delta t$ peak shape vs.\ rate, ADC
    (gain) spectra on- vs.\ off-spill, efficiency vs.\ instantaneous rate
    within the spill.
\end{enumerate}

\textbf{Hero plots:} \textbf{efficiency vs.\ trigger rate} (one point per
dataset, per geometry --- this is the money plot of Act~1); on-spill pad
occupancy map at 15~kHz; ADC MPV vs.\ rate (gain stability).

\subsection*{WP3 --- Pion runs \emph{(new analysis --- start early, highest uncertainty)}}

\begin{enumerate}[itemsep=1pt]
  \item Locate pion run numbers/conditions in the elog + run tables; add a
    \code{"pion"} entry to the \code{DATASETS} registries (the registry
    pattern means nothing else changes).
  \item Same chain as WP2: spill mask $\rightarrow$ occupancy map
    $\rightarrow$ coincidence efficiency, but now the key variable is
    \textbf{local rate density} (hits/pad/s in the beam-spot core vs.\ periphery).
  \item Efficiency and MPV as a function of local pad rate --- one curve
    combining muon + pion points extends the rate reach by an order of
    magnitude in a single figure.
\end{enumerate}

\textbf{Hero plots:} pion beam-spot occupancy map (localized spot on the pad
geometry --- visually striking); \textbf{efficiency (and/or MPV) vs.\ local pad
rate}, muon + pion points on one axis.

\subsection*{WP4 --- Geometry comparison \emph{(synthesis of WP1--3, cheap once they exist)}}

Large fan-pad vs.\ small pad detectors at identical $(sg, snt, \text{rate})$:
SNR, efficiency, occupancy per unit area. State the verdict explicitly: which
geometry carries forward, and why the new cosmic-bench detectors look the way
they do. This is the hinge between Act~1 and Act~2.

\textbf{Hero plot:} two-panel same-scale comparison (SNR map / efficiency map,
large vs.\ small).

\subsection*{WP5 --- Cosmic-bench characterization of the new detectors \emph{(pipeline exists --- needs det2 completion)}}

\begin{enumerate}[itemsep=1pt]
  \item Finish det2: HV scan (stage 11 equivalent for det2), sliding
    efficiency map, timing (commit stage 13), longer statistics run.
  \item Produce the \emph{presentation} versions: pad-tile efficiency map at
    best HV, HV turn-on curve det1 + det2 overlaid, TOA resolution plot.
  \item Stability narrative (optional but strong): duty-cycle/long-run
    monitoring plot showing how the bench catches gain instability --- frame
    as methodology, not as a failure.
\end{enumerate}

\textbf{Hero plots:} pad-tile efficiency map with M3 tracks (the ``we see
every pad'' plot); HV turn-on $\varepsilon(\mathrm{HV})$; timing residual
distribution with $\sigma \approx 46$~ns; one spark/stability monitoring time
series.

\subsection*{WP6 --- Cross-campaign synthesis \emph{(the conclusion slide --- do last, 2--3 days)}}

\begin{itemize}[itemsep=1pt]
  \item One table/figure: detector generation $\times$ electronics $\times$
    environment $\rightarrow$ efficiency, SNR, timing, rate reach. This is the
    single slide people will photograph.
  \item Explicit outlook: next beam test $=$ bench-validated det2 ($+$det1 if
    recovered) $+$ VMM at the WP1 optimal config $+$ rate scan informed by
    WP2/3.
\end{itemize}

%=====================================================================
\section{Suggested timeline (back-count from the conference)}
%=====================================================================

\begin{table}[h!]
\centering\small
\begin{tabular}{L{3.6cm} L{11.6cm}}
\toprule
\textbf{When} & \textbf{Milestone} \\
\midrule
Now $\rightarrow$ T$-$8 weeks &
  WP3 first pass (pion data located and read); WP2 efficiency-vs-rate sweep
  running; det2 HV scan taken on the bench \\
\addlinespace
T$-$8 $\rightarrow$ T$-$5 weeks &
  WP1 frozen; WP2/WP3 result plots v1; det2 characterization complete;
  commit stage 13 \\
\addlinespace
T$-$5 $\rightarrow$ T$-$3 weeks &
  WP4 + WP5 presentation plots finalized (consistent style, labels, fonts);
  WP6 synthesis figure \\
\addlinespace
T$-$3 $\rightarrow$ T$-$1 weeks &
  Slides assembled; internal dry run; only cosmetic plot changes after this \\
\addlinespace
T$-$1 week &
  Freeze. Backup slides: methodology (spill mask, $\Delta t$ fit, MAD noise),
  det1 stability study, mapping validation vs.\ M3 \\
\bottomrule
\end{tabular}
\end{table}

%=====================================================================
\section{Slide skeleton (\texorpdfstring{$\approx$}{~}15 min talk)}
%=====================================================================

\begin{enumerate}[itemsep=1pt]
  \item Motivation \& detector concept (P2 basket pad-Micromegas, target
    environment) --- 1--2 slides
  \item \textbf{Act 1}: SPS setup photo + two geometries + VMM3a readout --- 1 slide
  \item SNR optimization: scan matrix + best-config spectra --- 2 slides
  \item Rate performance: efficiency vs.\ rate (muons), pion beam-spot +
    efficiency vs.\ local rate --- 2 slides
  \item Geometry verdict --- 1 slide
  \item \textbf{Act 2}: cosmic bench + M3 telescope setup --- 1 slide
  \item New detectors: pad-tile efficiency map, HV turn-on, timing --- 2 slides
  \item Stability \& QA methodology (spark veto, pedestal, long-run
    monitoring) --- 1 slide
  \item \textbf{Act 3}: synthesis table + next beam test plan --- 1--2 slides
\end{enumerate}

%=====================================================================
\section{Practical notes}
%=====================================================================

\begin{itemize}[itemsep=1pt]
  \item Every plot destined for the talk should get a presentation-style pass
    (single consistent font/size, no debug titles, units on every axis); both
    pipelines already save PNG+PDF, so keep the PDF versions for slides.
  \item The \code{DATASETS} registry pattern (SNR, trigger, cosmic pipelines)
    is the mechanism for every new campaign entry --- pion runs, det2 HV scan
    --- one dict entry each, nothing else changes.
  \item Open data questions to resolve early: pion run list/conditions (elog),
    the empty 365~V m3 file (re-fetch from \code{rays\_daplxa}), whether det1
    recovers enough to appear as a second bench detector.
\end{itemize}

\end{document}
